% Part: computability
% Chapter: computability-theory
% Section: total

\documentclass[../../../include/open-logic-section]{subfiles}

\begin{document}

\olfileid[pt]{cmp}{thy}{tot}
\olsection{A totalidade é indecidível}

Consideremos mais um exemplo do uso do teorema $s$-$m$-$n$ para
mostrar que algo não é computável. Seja $\fn{Tot}$ o conjunto dos
índices das funções computáveis totais, isto é,
\[
\fn{Tot} = \Setabs{x}{\text{para todo $y$, $\cfind{x}(y)\fdefined$}}.
\]

\begin{prop}
\ollabel{prop:total}
$\fn{Tot}$ não é computável.
\end{prop}

\begin{proof}
Para ver que $\fn{Tot}$ não é computável, basta mostrar que $K$ é
redutível a ele. Seja $h(x,y)$ definida por
\[
h(x,y) \simeq
\begin{cases}
0 & \text{se $x \in K$} \\
\fundefined & \text{caso contrário}
\end{cases}
\]
Observe que $h(x,y)$ não depende de modo algum de $y$. Não deve ser
difícil ver que $h$~é parcialmente computável: nas entradas $x, y$,
calculamos~$h$ simulando primeiro a função~$\cfind{x}$ na entrada~$x$;
se essa computação parar, $h(x,y)$ produzirá $0$ como saída e parará.
Assim, $h(x,y)$ é simplesmente $\Zero(\umin{s}{T(x,x,s)})$, em que
$\Zero$ é a função constante zero.

Usando o teorema $s$-$m$-$n$, existe uma função recursiva
primitiva~$k(x)$ tal que, para todos $x$ e~$y$,
\[
\cfind{k(x)}(y) =
\begin{cases}
0 & \text{se $x \in K$} \\
\fundefined & \text{caso contrário}
\end{cases}
\]
Logo, $\cfind{k(x)}$ é total se $x \in K$ e é indefinida caso
contrário. Portanto, $k$ é uma redução de $K$ a~$\fn{Tot}$.
\end{proof}

\begin{digress}
Verifica-se que $\fn{Tot}$ nem sequer é computavelmente enumerável---sua
complexidade está mais acima na ``hierarquia aritmética''. Mas não
nos ocuparemos aqui desse fortalecimento.
\end{digress}

\end{document}
